\documentclass[12pt]{article}
% Shared print layout. The web's provisional design and MathJax are independent.
\usepackage[utf8]{inputenc}
\usepackage[T1]{fontenc}
\usepackage[spanish,es-nodecimaldot]{babel}
\usepackage{lmodern}
\usepackage{amsmath,amssymb}
\usepackage{graphicx}
\usepackage{geometry}
\geometry{a4paper,left=20mm,right=20mm,top=25mm,bottom=22mm,headheight=16pt,headsep=8mm}
\usepackage{fancyhdr,lastpage}
\usepackage{enumitem}
\usepackage{xcolor}
\usepackage{pgfplots}
\pgfplotsset{compat=1.18}
\usepgfplotslibrary{fillbetween}
\usepackage[hidelinks,pdfencoding=auto]{hyperref}
\hypersetup{pdflang=es-ES,pdfauthor={},pdfsubject={Matemáticas Aplicadas a las Ciencias Sociales II}}
\pagestyle{fancy}
\fancyhf{}
\fancyhead[L]{\small Matemáticas Aplicadas a las Ciencias Sociales II}
\fancyhead[R]{\small 2.\textsuperscript{o} Bachillerato}
\fancyfoot[L]{\footnotesize Sistemas de ecuaciones II | Edición revisada 1.0}
\fancyfoot[R]{\footnotesize Página \thepage\ de \pageref{LastPage}}
\renewcommand{\headrulewidth}{0.4pt}
\setlength{\parindent}{0pt}
\setlength{\parskip}{0.5em}
\setlist{itemsep=0.3em,topsep=0.3em}
\newcommand{\apartado}[1]{\subsection*{#1}}
\pgfplotsset{
  revision/.style={
    width=0.94\linewidth,height=7.5cm,
    axis lines=middle,grid=major,grid style={gray!18},
    tick label style={font=\small},label style={font=\small},
    xlabel={$x$},ylabel={$y$},samples=160,
    legend style={at={(0.5,-0.16)},anchor=north,draw=none,font=\small},
    clip=true
  }
}

\hypersetup{pdftitle={Sistemas de ecuaciones II - Examen revisado}}
\begin{document}
\begin{minipage}{0.17\linewidth}
\includegraphics[width=0.88\linewidth]{logo-al-andalus.png}
\end{minipage}\hfill
\begin{minipage}{0.79\linewidth}
{\Large\bfseries Sistemas de ecuaciones II}\\[0.4em]
Examen 2. Duración: \textbf{55 minutos}.\\
Puntuación total: \textbf{10 puntos}.
\end{minipage}

Apellidos y nombre: \hrulefill

Fecha: \hrulefill

Lee las preguntas y responde de forma razonada. Numera las páginas de tu respuesta. Se permite el uso de calculadora.
\section*{Ejercicio 1 \hfill 2 puntos}
\textit{Tiempo orientativo: 9 minutos.} Resuelve el sistema y encuentra $x$, $y$, $z$:
\[
\begin{pmatrix}2&1&-1\\0&3&5\\0&0&1\end{pmatrix}
\begin{pmatrix}x\\y\\z\end{pmatrix}
=\begin{pmatrix}4\\7\\1\end{pmatrix}.
\]
\textbf{Pista del original:} en ocasiones es más fácil escribirlo como sistema y reducirlo a solo dos variables si es posible.
\section*{Ejercicio 2 \hfill 4 puntos}
\textit{Tiempo orientativo: 23 minutos.} Una empresa de logística realiza envíos de paquetes pequeños, medianos y grandes. En una semana procesó 150 envíos y obtuvo ingresos de 2\,450 euros. El número de envíos de paquetes pequeños es el doble de la diferencia entre los de paquetes medianos y los de paquetes grandes, en ese orden.

Los precios por envío son: pequeño, 10 euros; mediano, 15 euros; grande, 20 euros. Determina el número de envíos de cada tipo.
\section*{Ejercicio 3 \hfill 4 puntos}
\textit{Tiempo orientativo: 23 minutos.} Dado el sistema con parámetro $\lambda$:
\[
\begin{cases}x+y+z=3,\\2x+3y+\lambda z=6,\\4x+5y+7z=10.\end{cases}
\]
\begin{enumerate}[label=(\alph*)]
\item (2,5 puntos) Discute el sistema en función de los valores de $\lambda$.
\item (1,5 puntos) Si es posible, resuelve el sistema para $\lambda=0$.
\end{enumerate}
\end{document}
